% =============================================================================
% Appendix B: Full Configuration Reference
% =============================================================================
\mychapter{Appendix B --- Full Configuration Reference}

This appendix provides the complete configuration parameter reference for the XCT defect detection
pipeline, as actually defined in \texttt{config.py} in the repository root as of August 2026. Three
values are corrected from the mid-term submission's Appendix B (marked below); all other values were
verified to match. Bernsen thresholding parameters, added for the PODFAM phase and not present at
mid-term, are listed separately at the end of this appendix and documented further in Part II,
Section~9.3.

\mysection{B.1 Preprocessing Parameters}
\begin{table}[htbp]
\centering
\small
\begin{tabular}{p{4.0cm}p{2.8cm}p{6.4cm}}
\toprule
\textbf{Parameter} & \textbf{Default} & \textbf{Description} \\
\midrule
\param{NORM\_LOW\_PERCENTILE} & 1 & Lower clip percentile \\
\param{NORM\_HIGH\_PERCENTILE} & 99 & Upper clip percentile \\
\param{NORM\_EPS} & $10^{-7}$ & Division stability constant \\
\param{MEDIAN\_KERNEL\_SIZE} & 3 & Pre-NLM median filter kernel size (not listed at mid-term) \\
\param{BHC\_POLY\_DEGREE} & 3 & BHC polynomial degree \\
\param{NLM\_PATCH\_SIZE} & 5 & NLM patch size (pixels) \\
\param{NLM\_PATCH\_DIST} & 6 & NLM search distance (pixels) \\
\param{NLM\_H\_FACTOR} & \textbf{0.6} (CORRECTED from 0.08) & NLM filter strength multiplier,
  $h = \texttt{NLM\_H\_FACTOR} \times \hat{\sigma}$, adaptive \\
\param{RING\_FILTER\_RADIUS} & 15 & Polar median filter radius \\
\bottomrule
\end{tabular}
\caption{Preprocessing configuration parameters. \texttt{NLM\_H\_FACTOR} corrected; see the Note on
  this combined document.}
\end{table}

\mysection{B.2 Label Generation Parameters (Otsu / morphological, Part I)}
\begin{table}[htbp]
\centering
\begin{tabular}{p{4.0cm}p{2.8cm}p{6.4cm}}
\toprule
\textbf{Parameter} & \textbf{Default} & \textbf{Description} \\
\midrule
\param{MORPH\_OPEN\_SIZE} & 3 & Morphological opening element size \\
\param{MORPH\_CLOSE\_SIZE} & 3 & Morphological closing element size \\
\param{MIN\_DEFECT\_SIZE} & 5 & Minimum component size (voxels) \\
\bottomrule
\end{tabular}
\caption{Automatic label generation configuration parameters (Otsu-based, Part I). All values
  verified to match the current codebase.}
\end{table}

\mysection{B.3 Dataset and Patch Parameters}
\begin{table}[htbp]
\centering
\begin{tabular}{p{4.0cm}p{2.8cm}p{6.4cm}}
\toprule
\textbf{Parameter} & \textbf{Default} & \textbf{Description} \\
\midrule
\param{PATCH\_SIZE} & 256 & Patch size in pixels \\
\param{PATCH\_STRIDE} & 128 & Sliding window stride (50\% overlap) \\
\param{FG\_BG\_RATIO} & (1, 3) & Foreground to background patch ratio \\
\param{MIN\_FG\_PIXELS} & 10 & Minimum defect pixels per foreground patch \\
\param{VAL\_SPLIT} & \textbf{0.2} (CORRECTED from 0.15) & Validation set fraction \\
\param{TEST\_SPLIT} & \textbf{0.1} (CORRECTED from 0.15) & Test set fraction \\
\bottomrule
\end{tabular}
\caption{Dataset and patch extraction configuration parameters. \texttt{VAL\_SPLIT}/\texttt{TEST\_SPLIT}
  corrected (actual split is 70/20/10, not 70/15/15); see Section~3.2.2 and the Note on this combined
  document.}
\end{table}

\mysection{B.4 Training Parameters}
\begin{table}[htbp]
\centering
\small
\begin{tabular}{p{4.0cm}p{2.8cm}p{6.4cm}}
\toprule
\textbf{Parameter} & \textbf{Default} & \textbf{Description} \\
\midrule
\param{BATCH\_SIZE} & 8 & Training batch size (2D) \\
\param{NUM\_EPOCHS} & 50 & Maximum training epochs \\
\param{LEARNING\_RATE} & $4\times10^{-5}$ & Adam optimiser learning rate \\
\param{WEIGHT\_DECAY} & $10^{-5}$ & Adam weight decay (not listed at mid-term) \\
\param{LOSS\_FUNCTION} & \texttt{dice\_focal} & Active loss function \\
\param{FOCAL\_ALPHA} & 0.25 & Focal loss alpha parameter \\
\param{FOCAL\_GAMMA} & 2.0 & Focal loss gamma parameter \\
\param{DICE\_FOCAL\_LAMBDA} & 0.5 & Dice weight in combined loss \\
\param{EARLY\_STOP\_PATIENCE} & 10 & Early stopping patience (epochs) \\
\param{SCHEDULER\_PATIENCE} & 5 & LR scheduler patience (not listed at mid-term) \\
\param{DEVICE} & \texttt{cuda} & Training device (\texttt{cuda} or \texttt{cpu}) \\
\param{CKPT\_DIR} & \textbf{\texttt{artifacts/}} (CORRECTED from \texttt{checkpoints/}) & Model
  checkpoint directory \\
\bottomrule
\end{tabular}
\caption{Model training configuration parameters. \texttt{CKPT\_DIR} corrected to match the actual
  output directory used by \url{training/trainer.py}.}
\end{table}

\mysection{B.5 Evaluation Parameters}
\begin{table}[htbp]
\centering
\begin{tabular}{p{4.0cm}p{2.8cm}p{6.4cm}}
\toprule
\textbf{Parameter} & \textbf{Default} & \textbf{Description} \\
\midrule
\param{ACCEPTANCE\_DICE} & 0.75 & Minimum Dice on held-out test set \\
\param{ACCEPTANCE\_IOU} & 0.60 & Minimum IoU on held-out test set \\
\param{ACCEPTANCE\_REC} & 0.80 & Minimum Recall for defects $> 50\,\mu$m \\
\param{DICE\_THRESHOLD} & 0.5 & Sigmoid output threshold for binary mask \\
\bottomrule
\end{tabular}
\caption{Evaluation and acceptance criteria configuration parameters. All values verified to match
  the current codebase.}
\end{table}

\mysection{B.6 Augmentation Parameters}
\begin{table}[htbp]
\centering
\small
\begin{tabular}{p{4.0cm}p{2.8cm}p{6.4cm}}
\toprule
\textbf{Parameter} & \textbf{Default} & \textbf{Description} \\
\midrule
\param{AUG\_FLIP\_PROB} & 0.5 & Flip probability \\
\param{AUG\_ROTATE\_PROB} & 0.5 & Rotation probability \\
\param{AUG\_ELASTIC\_PROB} & 0.3 & Elastic deformation probability \\
\param{AUG\_ELASTIC\_ALPHA} & 34 & Elastic deformation magnitude \\
\param{AUG\_ELASTIC\_SIGMA} & 4 & Elastic deformation smoothness \\
\param{AUG\_INTENSITY\_PROB} & 0.5 & Intensity scaling probability \\
\param{AUG\_INTENSITY\_RANGE} & (0.9, 1.1) & Intensity scaling range \\
\param{AUG\_NOISE\_PROB} & 0.5 & Gaussian noise probability \\
\param{AUG\_NOISE\_STD\_RANGE} & (0.01, 0.05) & Noise standard deviation range \\
\param{AUG\_GAMMA\_PROB} & 0.5 & Gamma correction probability \\
\param{AUG\_GAMMA\_RANGE} & (0.8, 1.2) & Gamma correction range \\
\bottomrule
\end{tabular}
\caption{Data augmentation configuration parameters. All values verified to match the current
  codebase.}
\end{table}

\mysection{B.7 Bernsen Thresholding Parameters (PODFAM phase, Part II --- added since mid-term)}
\begin{table}[htbp]
\centering
\small
\begin{tabular}{p{4.6cm}p{2.4cm}p{6.2cm}}
\toprule
\textbf{Parameter} & \textbf{Default} & \textbf{Description} \\
\midrule
\param{BERNSEN\_RADIUS} & 5 & Local window radius in pixels (Kim et al.\ 2017) \\
\param{BERNSEN\_LOW\_CONTRAST\_THRESH} & 128 & Fixed threshold for low-contrast regions \\
\param{BERNSEN\_DCT\_AUTO} & \texttt{True} & Compute DCT per image (recommended) \\
\param{BERNSEN\_DCT\_STD\_MULTIPLIER} & 18 & From Kim et al.\ 2017 (18$\times$ average std) \\
\param{BERNSEN\_DCT\_N\_LOCATIONS} & 5 & Grid points per axis for std sampling \\
\param{BERNSEN\_DCT\_WINDOW\_RADIUS} & 5 & Local window radius for std measurement \\
\param{BERNSEN\_DCT\_MIN} & 5 & Minimum DCT floor (prevents under-segmentation) \\
\param{BERNSEN\_DCT} & 30 & Fixed fallback DCT if \texttt{BERNSEN\_DCT\_AUTO} is \texttt{False} \\
\param{PIXEL\_SIZE\_UM} & 1.0 & Uncalibrated placeholder; see Part II Section~12 \\
\param{USE\_SAMPLE\_MASK} & \texttt{True} & Apply circular sample-boundary mask before
  thresholding \\
\param{SAMPLE\_MASK\_EROSION\_RADIUS} & 5 & Erosion margin below detected sample surface \\
\bottomrule
\end{tabular}
\caption{Bernsen thresholding configuration parameters, not present at mid-term.}
\end{table}
